% =====================================================================
%  Lecture 18 -- 5G Service-Based Architecture (SBA)      (English edition)
% =====================================================================
\documentclass[11pt,a4paper]{article}
\providecommand\ACNroot{../..}
\input{\ACNroot/templates/lecture-en}

\ACNset{number}{18}
\ACNset{title}{5G Service-Based Architecture (SBA)}
\ACNset{shorttitle}{5G SBA}
\ACNset{titleother}{معماری سرویس‌محور نسل پنجم (SBA)}
\ACNset{part}{Part III -- Software-driven architectures (Lectures 15--18)}
\ACNset{time}{6--8 hours (3 h reading, 3--5 h exercises)}
\ACNset{level}{Intermediate}
\ACNset{prereq}{Lectures 10, 14 and 16; HTTP, REST and JSON basics}
\ACNset{gen}{5}

\begin{document}
\ACNcover

% ---------------------------------------------------------------------
\section{Learning objectives}
After studying this lecture you should be able to:
\begin{objectives}
  \item \textbf{Explain} the principles of the service-based architecture: NF services, producers and consumers,
        request/response and subscribe/notify.
  \item \textbf{Describe} the protocol stack of service-based interfaces (HTTP/2, TLS, JSON, OpenAPI) and the naming
        and URI conventions of NF services.
  \item \textbf{Trace} NF registration, discovery, authorisation and service invocation through the NRF.
  \item \textbf{Compare} direct and indirect (SCP) communication and \textbf{explain} the SEPP's role in roaming.
  \item \textbf{Quantify} connection needs, discovery load and latency of service-based signalling.
  \item \textbf{Contrast} SBA with the point-to-point Diameter/GTP-C design of the EPC.
\end{objectives}

% ---------------------------------------------------------------------
\section{Prerequisite knowledge}
You need the Diameter-based EPC interfaces (Lecture~10) as a contrast, the 5G overview of Lecture~14, the
cloud-native principles of Lecture~16, and familiarity with HTTP methods (GET, POST, PUT, PATCH, DELETE), status codes,
URIs and JSON.

\paragraph{Key terms} The vocabulary of this lecture with its Persian equivalents:
\begin{keyterms}
SBA / SBI & \fa{معماری / واسط سرویس‌محور} & Architecture in which NFs offer services over a common interface. \\
NF service & \fa{خدمت کارکرد شبکه} & Capability exposed by an NF (e.g.\ Nsmf\_PDUSession). \\
Producer / consumer & \fa{تولیدکننده / مصرف‌کننده} & NF offering / using a service. \\
NRF & \fa{تابع مخزن شبکه} & Registry for NF registration, discovery and access tokens. \\
SCP & \fa{پراکسی ارتباط خدمت} & Proxy for indirect communication and delegated discovery. \\
SEPP & \fa{پراکسی حفاظت لبهٔ امنیت} & Edge proxy securing SBI signalling between PLMNs (N32). \\
OpenAPI & \fa{OpenAPI} & Machine-readable description of a REST API. \\
HTTP/2 stream & \fa{جریان HTTP/2} & One request/response exchange multiplexed on a connection. \\
OAuth 2.0 & \fa{OAuth 2.0} & Authorisation framework; the NRF issues access tokens \cite{rfc6749}. \\
NF profile & \fa{نمایهٔ NF} & Description of an NF instance stored in the NRF. \\
\end{keyterms}

% ---------------------------------------------------------------------
\section{Introduction and motivation}
The EPC was designed as a set of boxes connected by point-to-point interfaces: S6a between MME and HSS, S11 between MME
and S-GW, Gx between P-GW and PCRF, each with its own protocol (Diameter, GTPv2-C) and its own peer configuration.
Adding a function or a consumer meant defining a new interface and configuring every peer. In a cloud-native core, where
instances come and go with load, this rigidity becomes the bottleneck.

The 5G core adopts the patterns of web-scale software: each network function exposes \emph{services} through a common
\emph{service-based interface} (SBI) built on HTTP/2 and JSON; any authorised NF can consume them; NFs register with a
repository function (NRF) and discover each other at run time \cite{ts23501,ts29500}. This lecture closes Part~III by
showing how SDN, NFV and web technology come together in the 5G control plane; Part~IV then uses these services to build
the 5G procedures.

% ---------------------------------------------------------------------
\section{Fundamental concepts}

\subsection{Network functions and services}
The control plane of the 5G core consists of NFs connected to the SBI (\cref{fig:sba}); the user plane (UPF) and the
access network keep classical reference points (N2, N3, N4, N6). \Cref{tab:nfs} lists the main NFs.

\begin{table}[htbp]
\centering\small
\caption{Main 5G core network functions and example services \cite{ts23501}.}\label{tab:nfs}
\rowcolors{2}{acnmist}{white}
\begin{tabularx}{\linewidth}{@{}l Y l@{}}
\toprule
\rowcolor{white}NF & Role & Example service \\
\midrule
AMF  & Access and mobility management, NAS termination      & Namf\_Communication \\
SMF  & Session management, UPF control                      & Nsmf\_PDUSession \\
UPF  & User-plane forwarding (not on SBI)                   & -- (N4/PFCP) \\
AUSF & Authentication server                                & Nausf\_UEAuthentication \\
UDM / UDR & Subscription data management / repository       & Nudm\_SDM, Nudm\_UECM \\
PCF  & Policy control                                       & Npcf\_SMPolicyControl \\
NRF  & NF repository: registration, discovery, tokens       & Nnrf\_NFDiscovery \\
NSSF & Network slice selection                              & Nnssf\_NSSelection \\
NEF  & Exposure of network capabilities to external AFs     & Nnef\_EventExposure \\
CHF  & Charging                                             & Nchf\_ConvergedCharging \\
NWDAF & Network data analytics                              & Nnwdaf\_AnalyticsInfo \\
SCP, SEPP & Indirect communication, inter-PLMN security     & -- (proxies) \\
\bottomrule
\end{tabularx}
\end{table}

\begin{acnfigure}
  \resizebox{\linewidth}{!}{\acnfig{lecture-18/sba}}
  \caption{5G system with service-based control plane (SBI) and reference points to the RAN and UPF.}
  \label{fig:sba}
\end{acnfigure}

\subsection{Producers, consumers and interaction patterns}
A \emph{producer} NF offers a service; a \emph{consumer} NF uses it. Services are named
\emph{N\textless nf\textgreater\_\allowbreak\textless Service\textgreater\_\allowbreak\textless Operation\textgreater}, e.g.\ Nsmf\_\allowbreak PDUSession\_\allowbreak CreateSMContext.
Two interaction patterns exist:
\begin{itemize}
  \item \textbf{request/response}: the consumer requests an operation and receives the result;
  \item \textbf{subscribe/notify}: the consumer subscribes to events (e.g.\ AMF subscribes to UDM for subscription changes)
        and the producer sends notifications to a callback URI given by the consumer.
\end{itemize}
Services are decoupled from the NFs that use them: a new consumer can use an existing service without changing the producer
-- the main architectural gain over point-to-point interfaces.

\subsection{Protocol stack}
SBIs use HTTP/2 \cite{rfc9113} over TCP with TLS, JSON bodies, and APIs described in OpenAPI 3.0 documents published with the
specifications \cite{ts29500,ts29501}. Resources are identified by URIs of the form
\nolinkurl{{apiRoot}/<api-name>/<version>/<resource>}, e.g.\
\nolinkurl{{apiRoot}/nsmf-pdusession/v1/sm-contexts}. HTTP methods map to operations (\cref{tab:http}); errors are returned
as HTTP status codes with a JSON problem description.

\begin{table}[htbp]
\centering\small
\caption{Use of HTTP methods on service-based interfaces \cite{ts29501}.}\label{tab:http}
\rowcolors{2}{acnmist}{white}
\begin{tabularx}{\linewidth}{@{}l Y l@{}}
\toprule
\rowcolor{white}Method & Meaning & Typical answer \\
\midrule
POST   & Create a resource in a collection, or invoke a custom operation & 201 Created (Location), 200 OK \\
GET    & Read a resource or query a collection                  & 200 OK \\
PUT    & Create or replace a resource with a client-chosen URI  & 201 Created, 200/204 \\
PATCH  & Partially modify a resource (JSON Patch / merge patch) & 200 OK, 204 No Content \\
DELETE & Remove a resource                                      & 204 No Content \\
\bottomrule
\end{tabularx}
\end{table}

% ---------------------------------------------------------------------
\section{Detailed technical explanation}

\subsection{The NRF}
The NRF provides three services \cite{ts29510}:
\begin{itemize}
  \item \textbf{Nnrf\_NFManagement}: an NF instance registers its \emph{NF profile} (instance ID, type, PLMN, supported slices,
        services and their endpoints, capacity, priority) with PUT, keeps it alive with periodic heartbeats, updates it with
        PATCH and deregisters with DELETE; NFs can subscribe to changes of other profiles.
  \item \textbf{Nnrf\_NFDiscovery}: a consumer queries for producers by target NF type and criteria (slice, DNN, TAI,
        SUPI range, \ldots); the answer contains matching profiles and a validity period during which they may be cached.
  \item \textbf{Nnrf\_AccessToken}: the NRF acts as OAuth~2.0 authorisation server and issues access tokens (client
        credentials grant) that a producer verifies before serving a request \cite{rfc6749}.
\end{itemize}

\subsection{Direct and indirect communication}
In \emph{direct} communication the consumer sends requests straight to the producer, either from configuration (option A)
or after NRF discovery (option B). In \emph{indirect} communication (Release~16) requests pass through a \emph{service
communication proxy} (SCP): the consumer either discovers the producer itself (option C) or delegates discovery and selection
to the SCP (option D) (\cref{fig:comm}). The SCP plays a role similar to a Diameter routing agent or a service mesh: load
balancing, routing, overload handling and observability in one place, at the cost of an extra hop.

\begin{acnfigure}
  \resizebox{0.9\linewidth}{!}{\acnfig{lecture-18/comm-options}}
  \caption{Direct communication (options A/B) and indirect communication via an SCP (options C/D).}
  \label{fig:comm}
\end{acnfigure}

\subsection{Resilience: NF sets and binding}
NF instances that share context (e.g.\ via a common data store) can form an \emph{NF set}; any instance of the set can serve
a UE's requests, so a failed instance is replaced without losing the session. \emph{Binding indications} tell consumers and
SCPs which instance or set should receive subsequent requests or notifications for a given context. This is the SBA
counterpart of the stateless-instance design of cloud-native functions (Lecture~16).

\subsection{Roaming and the SEPP}
Between PLMNs, NFs never talk directly: each PLMN has a \emph{security edge protection proxy} (SEPP) at its border, and the
SEPPs interconnect over N32. The SEPPs negotiate security on N32-c and protect forwarded messages on N32-f, either with TLS
end to end between the SEPPs or with application-layer protection when intermediate IPX providers must be able to read or
modify selected fields \cite{ts33501}. The SEPP thus plays the role of the Diameter edge agent of the EPC, with stronger
security.

\subsection{Service operations in practice}
A few examples of resource-style operations (simplified) \cite{ts29510}:
\begin{itemize}
  \item registration: \nolinkurl{PUT {apiRoot}/nnrf-nfm/v1/nf-instances/{nfInstanceId}} with the NF profile;
  \item discovery: \nolinkurl{GET {apiRoot}/nnrf-disc/v1/nf-instances?target-nf-type=SMF&requester-nf-type=AMF};
  \item session creation: \nolinkurl{POST {apiRoot}/nsmf-pdusession/v1/sm-contexts} with a JSON body containing SUPI, DNN, S-NSSAI
        and the N1 message; the SMF answers 201 with the URI of the new SM context.
\end{itemize}

% ---------------------------------------------------------------------
\section{Mathematical models}

\subsection{Concurrency on HTTP/2 connections}\label{sec:inflight}
\paragraph{Model} A consumer sends $\lambda$ requests per second to a producer; each request is outstanding for a response time
$R$ (network round trip plus processing). By Little's law the number of concurrent requests is
\begin{equation}\label{eq:inflight}
  N_{\mathrm{inflight}} = \lambda R .
\end{equation}
With HTTP/1.1 without pipelining, each outstanding request needs its own connection; with HTTP/2 a single connection carries up to
the peer's maximum number of concurrent streams (recommended to be at least 100 \cite{rfc9113}).
\paragraph{Example} $\lambda = 10^4$\,s$^{-1}$, $R = 5$\,ms: $N_{\mathrm{inflight}} = 50$ -- fifty HTTP/1.1 connections or one HTTP/2
connection (in practice a few, for redundancy and load spreading).

\subsection{Discovery load and caching}\label{sec:cache}
\paragraph{Model} Without caching, every request triggers a discovery and the NRF sees the full request rate $\lambda$. If each of $C$
consumers caches discovery results for $K$ target types for a validity period $T$, the NRF sees approximately
\begin{equation}\label{eq:cache}
  \lambda_{\mathrm{NRF}} \approx \frac{C\,K}{T}.
\end{equation}
\paragraph{Example} $C = 100$, $K = 5$, $T = 60$\,s: 8.3 discoveries/s instead of $10^4$/s (\cref{fig:nrfload}).
\paragraph{Interpretation} Caching makes the NRF a light-weight directory, but cached profiles become stale: a consumer may keep
sending requests to an overloaded or failed instance until the cache expires or a notification arrives. The validity period
trades NRF load against reaction time.

\begin{acnfigure}
  \acnfig{lecture-18/nrf-load}
  \caption{NRF discovery load versus cache validity (\cref{eq:cache}).}
  \label{fig:nrfload}
\end{acnfigure}

\subsection{Latency of indirect communication}\label{sec:indirect}
\paragraph{Model} Let $T_{\mathrm{dir}}$ be the response time of a direct request. Routing through an SCP adds two network hops of
one-way delay $d$ each direction and SCP processing $t_{\mathrm{s}}$; with delegated discovery the SCP occasionally (probability
$p_{\mathrm{m}}$ of a cache miss) queries the NRF, costing $T_{\mathrm{NRF}}$:
\begin{equation}\label{eq:indirect}
  T_{\mathrm{ind}} \approx T_{\mathrm{dir}} + 2d + t_{\mathrm{s}} + p_{\mathrm{m}}\,T_{\mathrm{NRF}} .
\end{equation}
\paragraph{Example} $T_{\mathrm{dir}} = 3$\,ms, $d = 0.25$\,ms, $t_{\mathrm{s}} = 0.5$\,ms, $p_{\mathrm{m}} = 0.05$, $T_{\mathrm{NRF}} = 2$\,ms:
$T_{\mathrm{ind}} = 4.1$\,ms. A registration that needs ten sequential SBI exchanges gains about 11\,ms.

\subsection{Header compression}\label{sec:hpack}
\paragraph{Model} HTTP/2 compresses headers with HPACK: the first request on a connection sends headers of $H_1$ bytes; later
requests mostly reference table entries and cost $h$ bytes. The mean header size over $n$ requests is
\begin{equation}\label{eq:hpack}
  \bar H_n = \frac{H_1 + (n-1)\,h}{n} \;\xrightarrow{\;n\to\infty\;}\; h .
\end{equation}
\paragraph{Example} $H_1 = 400$\,bytes, $h = 20$\,bytes, $n = 100$: $\bar H = 23.8$\,bytes per request -- a reason to keep long-lived
connections between NFs.

% ---------------------------------------------------------------------
\section{Network architecture and diagrams}
In a deployment, NF instances run as containers in one or more data centres; each NF type is replicated and registered in the
NRF (itself redundant); SCPs or a service mesh connect them; SEPPs sit at the border; AMFs and UPFs connect to the RAN over N2 and
N3. Unlike the EPC, there is no fixed peer table: who talks to whom is decided at run time from the NRF's view, the slice and the
location of the UE (\cref{fig:sba}).

% ---------------------------------------------------------------------
\section{Step-by-step procedures}

\begin{procedurebox}{NF registration and discovery}\label{proc:register}
\begin{steps}
  \item A new SMF instance starts and sends PUT with its NF profile to the NRF (Nnrf\_NFManagement\_NFRegister); the NRF stores it
        and answers 201 with the heartbeat interval.
  \item The SMF sends periodic heartbeats; if they stop, the NRF marks it suspended and notifies subscribers.
  \item An AMF needing an SMF for a PDU session sends a discovery GET with target type SMF, DNN and S-NSSAI.
  \item The NRF returns matching profiles with a validity period; the AMF selects one (priority, capacity, locality) and caches the
        result.
\end{steps}
\end{procedurebox}

\begin{procedurebox}{Authorised service invocation}\label{proc:invoke}
\begin{steps}
  \item The AMF requests an access token from the NRF for scope ``nsmf-pdusession'' (Nnrf\_AccessToken).
  \item The NRF authenticates the AMF (TLS client certificate) and returns a signed token.
  \item The AMF sends POST to the SMF's sm-contexts resource with the token in the Authorization header.
  \item The SMF verifies the token's signature, scope and audience, creates the SM context and answers 201 with its URI; later
        operations use that URI.
\end{steps}
\end{procedurebox}

\begin{acnfigure}
  \resizebox{\linewidth}{!}{\acnfig{lecture-18/nrf-flow}}
  \caption{Registration, discovery, token and service invocation (\cref{proc:register,proc:invoke}).}
  \label{fig:nrf}
\end{acnfigure}

% ---------------------------------------------------------------------
\section{Worked examples}

\begin{examplebox}{Connections between NFs}\label{ex:inflight}
\textbf{Problem.} An AMF sends $2\times10^4$ requests/s to a set of SMFs with mean response time 8\,ms. How many requests are
outstanding, and how many HTTP/2 connections are needed if each allows 100 concurrent streams and connections should be at most
50\,\% utilised?

\textbf{Solution.} $N = 2\times10^4\times0.008 = 160$. With 50 streams used per connection, at least $\lceil160/50\rceil = 4$ connections,
spread over SMF instances for redundancy.
\end{examplebox}

\begin{examplebox}{NRF sizing}\label{ex:cache}
\textbf{Problem.} A core has 1000 consumer instances, each discovering 5 target types with a validity of 300\,s, plus 500 NF instances
sending a heartbeat every 10\,s. What load does the NRF see?

\textbf{Solution.} Discovery $1000\times5/300 = 16.7$\,s$^{-1}$; heartbeats $500/10 = 50$\,s$^{-1}$; total about 67 requests/s --
small, but critical: the NRF must be highly available, because new sessions cannot be placed without it once caches expire.
\end{examplebox}

\begin{examplebox}{SCP or direct?}\label{ex:scp}
\textbf{Problem.} With the numbers of \cref{sec:indirect}, by how much does indirect communication lengthen a procedure with 6
sequential SBI exchanges? What does the operator gain?

\textbf{Solution.} $6\times1.1 = 6.6$\,ms. The gains: consumers need no discovery logic or per-producer connection management; load
balancing, retries, overload control and monitoring are centralised; producers can be scaled or moved transparently.
\end{examplebox}

% ---------------------------------------------------------------------
\section{Technology comparison}
\begin{table}[htbp]
\centering\small
\caption{EPC point-to-point signalling versus 5G SBA.}\label{tab:cmp}
\rowcolors{2}{acnmist}{white}
\begin{tabularx}{\linewidth}{@{}l Y Y@{}}
\toprule
\rowcolor{white} & EPC (Diameter, GTPv2-C) & 5GC SBA \\
\midrule
Interface model   & Point-to-point reference points & Services on a common bus; any authorised consumer \\
Transport         & SCTP/TCP (Diameter), UDP (GTP-C) & TCP + TLS, HTTP/2 \\
Encoding          & Binary AVPs / IEs               & JSON, OpenAPI-defined \\
Peer discovery    & Static configuration, DNS, DRAs  & NRF registration and discovery \\
Interaction       & Request/answer                  & Request/response, subscribe/notify \\
Routing proxy     & DRA                             & SCP (optional) \\
Roaming edge      & DEA (Diameter), GTP firewalls   & SEPP with N32 security \\
Authorisation     & Implicit (peer configured)      & OAuth 2.0 tokens from the NRF \\
Extensibility     & New AVPs/IEs, new applications  & New services/operations, API versions \\
\bottomrule
\end{tabularx}
\end{table}

\paragraph{Not everything is service-based} N1 (NAS), N2 (NGAP), N4 (PFCP) and the user plane (N3, N9 with GTP-U) keep
point-to-point protocols, because they involve the RAN, the UE or high-rate forwarding where HTTP brings no benefit.

% ---------------------------------------------------------------------
\section{Key takeaways}
\begin{takeawaybox}
\begin{itemize}
  \item 5GC control-plane NFs expose services (Nnf\_Service\_Operation) on a common SBI: HTTP/2, TLS, JSON, OpenAPI.
  \item Request/response and subscribe/notify; producers are decoupled from consumers.
  \item The NRF handles registration, discovery and OAuth 2.0 tokens; SCP enables indirect communication; SEPP secures roaming.
  \item Little's law sizes concurrency ($\lambda R$); caching keeps NRF load at $CK/T$; an SCP adds about $2d + t_{\mathrm{s}}$ per exchange.
  \item N1, N2, N4 and the user plane remain point-to-point.
\end{itemize}
\end{takeawaybox}

% ---------------------------------------------------------------------
\section{Exercises}

\begin{exercisebox}{1}{Producers and consumers}\label{exr:pc}
During PDU session establishment the AMF calls Nsmf\_PDUSession\_CreateSMContext and the SMF calls Nudm\_SDM\_Get and
Npcf\_SMPolicyControl\_Create. Identify producer and consumer in each case.
\end{exercisebox}

\begin{exercisebox}{1}{Which method?}\label{exr:method}
Choose the HTTP method for: (a) reading a UE's subscribed slices; (b) creating an SM context; (c) registering an NF profile with
its known instance ID; (d) changing the load value in an NF profile; (e) releasing a policy association.
\end{exercisebox}

\begin{exercisebox}{2}{Connections}\label{exr:conn}
A UDM receives $5\times10^4$ requests/s from 40 consumer instances, response time 4\,ms. Compute the concurrency and compare the
number of connections needed with HTTP/1.1 (no pipelining) and HTTP/2 (100 streams per connection, one connection per consumer
instance). Is one HTTP/2 connection per consumer enough?
\end{exercisebox}

\begin{exercisebox}{2}{Cache validity}\label{exr:validity}
An operator wants the NRF discovery load below 50\,s$^{-1}$ with 2000 consumers and 6 target types. What is the minimum validity
period? What is the drawback of choosing it much larger?
\end{exercisebox}

\begin{exercisebox}{3}{Latency budget}\label{exr:budget}
A registration procedure has 12 sequential SBI exchanges with $T_{\mathrm{dir}} = 2.5$\,ms. The control-plane budget for registration in
the core is 50\,ms. Using \cref{eq:indirect} with $d = 0.5$\,ms, $t_{\mathrm{s}} = 0.8$\,ms, $p_{\mathrm{m}} = 0.1$, $T_{\mathrm{NRF}} = 3$\,ms, check
whether indirect communication fits. What changes if NFs are in two data centres 300\,km apart?
\end{exercisebox}

\begin{exercisebox}{3}{Design critique}\label{exr:critique}
Give two arguments for and two against using HTTP/2 + JSON instead of a binary protocol such as Diameter for core signalling, and
explain why N2 and N4 were not made service-based.
\end{exercisebox}

\subsection*{Solutions}
\addcontentsline{toc}{subsection}{Solutions}

\begin{solutionbox}{\ref{exr:pc}}
CreateSMContext: SMF producer, AMF consumer. SDM\_Get: UDM producer, SMF consumer. SMPolicyControl\_Create: PCF producer, SMF consumer.
\end{solutionbox}

\begin{solutionbox}{\ref{exr:method}}
(a) GET; (b) POST (to the sm-contexts collection); (c) PUT (to nf-instances/\{id\}); (d) PATCH; (e) DELETE (or POST to a custom
``delete'' operation where the API defines one).
\end{solutionbox}

\begin{solutionbox}{\ref{exr:conn}}
$N = 5\times10^4\times0.004 = 200$ concurrent requests. HTTP/1.1: 200 connections. HTTP/2: per consumer $200/40 = 5$ concurrent streams --
one connection per consumer suffices in capacity, but two are advisable for redundancy and for spreading load over UDM instances.
\end{solutionbox}

\begin{solutionbox}{\ref{exr:validity}}
$CK/T \le 50 \Rightarrow T \ge 2000\times6/50 = 240$\,s. A much larger validity makes consumers slow to notice new or removed producers;
they rely on NRF notifications or error responses to update, and load balancing over newly scaled instances is delayed.
\end{solutionbox}

\begin{solutionbox}{\ref{exr:budget}}
Per exchange: $2.5 + 1.0 + 0.8 + 0.3 = 4.6$\,ms; 12 exchanges: 55.2\,ms -- over the 50\,ms budget (direct: 30\,ms). With 300\,km between
data centres (one-way about 1.5\,ms in fibre), every cross-site exchange adds about 3\,ms of round trip: the procedure must keep
the NFs of one UE in one site (locality-aware discovery) or parallelise exchanges.
\end{solutionbox}

\begin{solutionbox}{\ref{exr:critique}}
For: widely available tools, libraries and skills; easy extensibility and versioning with OpenAPI; native fit with cloud
platforms (load balancers, service meshes, observability). Against: larger messages and parsing cost than binary TLV; TCP and TLS
connection management; more freedom (optional fields, versions) that can hurt interoperability. N2 and N4 involve the RAN and the
user plane, which need compact, low-latency, often hardware-assisted protocols and long-lived associations; their peers are not
generic consumers, so the SBA's decoupling brings little benefit there.
\end{solutionbox}

% ---------------------------------------------------------------------
\section{Further reading}
The SBA is defined in TS~23.501 \cite{ts23501} and its procedures in TS~23.502 \cite{ts23502}; TS~29.500 \cite{ts29500} and TS~29.501
\cite{ts29501} specify the technical realisation and API design rules; TS~29.510 \cite{ts29510} specifies the NRF; TS~33.501
\cite{ts33501} covers SBA security including the SEPP. HTTP/2 is RFC~9113 \cite{rfc9113} and OAuth~2.0 RFC~6749 \cite{rfc6749}.
Rommer et al.\ \cite{rommer2020core} give a textbook treatment.

\printbibliography[heading=subbibliography,title={\ACNstrRefs}]
\end{document}
